% Flavor catalog per-process draft.
% process_id: B026
% family: beauty
% owner_agent_id: PKA-B026
% writer_agent_id: null
% checker_agent_id: null
% opus_signoff_id: null
% Canonical metadata sidecar: B026.yaml
% Status is controlled by the YAML sidecar status_history, not this comment.

\section{B026: \(R_{D^*}\) in \(B\to D^*\tau\nu\)}

\subsection*{Process}
\textbf{Standard notation:} \(R_{D^*}\equiv
\mathcal{B}(B\to D^*\tau\nu_\tau)/
\mathcal{B}(B\to D^*\ell\nu_\ell)\), with \(\ell=e,\mu\).  This
beauty-family entry covers charged-current lepton-flavor-universality tests in
exclusive \(b\to c\ell\nu\) transitions with a vector charm meson in the final
state.

\subsection*{PDG or equivalent value}
The canonical current equivalent value is the HFLAV CKM 2025 joint
\(R_D\)-\(R_{D^*}\) average, updated on 2025-09-28 and snapshotted in
\texttt{flavor\_catalog/references/B026/hflav\_ckm2025\_rdrds.txt} with source
key \texttt{hflav\_ckm2025\_rdrds}:
\[
  R_{D^*}=0.281\pm0.011 .
\]
The same HFLAV table gives \(R_D=0.358\pm0.024\), correlation
\(\rho(R_D,R_{D^*})=-0.374\), and
\(\chi^2/\mathrm{dof}=16.683/14\) for the experimental combination.  HFLAV's
displayed arithmetic SM reference is \(R_D=0.296\pm0.004\) and
\(R_{D^*}=0.254\pm0.005\).  With the quoted covariance, HFLAV reports the
current combined \(R_D\)-\(R_{D^*}\) discrepancy as \(3.8\sigma\)
\((p=1.48\times10^{-4})\); the \(R_{D^*}\)-only pull against that reference is
\(2.3\sigma\).  All numerical values in this paragraph use source key
\texttt{hflav\_ckm2025\_rdrds}.

\subsection*{Relevance to RS / anarchic flavor}
\(R_{D^*}\) is not a neutral-meson \(\Delta F=2\) bound and does not constrain
the current quark-scan backend directly.  It becomes relevant in warped or
anarchic extensions with non-universal charged currents, third-generation
enhanced lepton couplings, charged-scalar exchange, \(W'\)-like vectors, or
leptoquark-motivated matching.  Because the observable is a ratio, many CKM and
normalization uncertainties cancel, making it a sensitive diagnostic once a
model predicts the \(b\to c\tau\nu\) Wilson coefficients.

\subsection*{Post-2008 developments}
The 2008 CFW/Perez-Randall-era baseline
(\texttt{cfw2008\_arxiv0804\_1954}) predates the \(R_D^{(*)}\) anomaly
program.  The post-2008 experimental sequence starts with BaBar's 2012 and
2013 excess in \(B\to D^{(*)}\tau\nu\)
(\texttt{babar2012\_arxiv1205\_5442};
\texttt{babar2013\_arxiv1303\_0571}), followed by Belle hadronic-tag,
tau-polarization, and semileptonic-tag measurements, LHCb muonic-tau and
three-prong-tau measurements, and Belle II inputs recorded in the HFLAV CKM
2025 average.  HFLAV now performs a correlated average rather than comparing
one experiment to one SM number (\texttt{hflav\_ckm2025\_rdrds};
\texttt{hflav\_ckm2025\_logfile}).  On the SM side,
\(B\to D^*\ell\nu\) form-factor theory changed materially after 2008: HFLAV
cites benchmark, FNAL/MILC, HPQCD, JLQCD, and FLAG inputs, with the local FLAG
snapshot stored as \texttt{flag2024\_arxiv2411\_04268}.  The residual tension
is therefore tied to correlated experimental averaging and form-factor
conventions, not to a missing basic SM prediction.

\subsection*{Constraint validity / model dependence}
The HFLAV average is a robust experimental summary, but interpreting it as an
RS constraint is highly model-dependent.  A live likelihood needs the
\((\bar c\,\Gamma b)(\bar\tau\,\Gamma\nu)\) charged-current operator basis,
neutrino-flavor assumptions, and the integrated \(B\to D^*\) response to
vector, scalar, pseudoscalar, and tensor operators.  A standalone Gaussian for
\(R_{D^*}\) would discard the documented correlation with \(R_D\); any serious
implementation should use the two-dimensional HFLAV covariance or a dedicated
likelihood.

\subsection*{Code coverage in this repo}
\textbf{NO.}  The PKA reran the required greps over
\texttt{quarkConstraints/}, \texttt{qcd/}, \texttt{flavorConstraints/},
\texttt{neutrinos/}, \texttt{yukawa/}, \texttt{warpConfig/},
\texttt{solvers/}, \texttt{scanParams/}, and \texttt{tests/}.  The broad
searches recover existing neutral-meson \(\Delta F=2\) and
\(\mu\to e\gamma\) code, but targeted searches for \(R(D^*)\), \(D^*\tau\),
\(B\to D^*\), \(b\to c\), charged-current, charged-Higgs, and leptoquark
observable code returned no implementation hits.

\subsection*{Implementation difficulty}
\textbf{HIGH.}  Promotion to a live constraint requires new charged-current
semileptonic matching, \(B\to D^*\) form-factor integration, and a correlated
\(R_D\)-\(R_{D^*}\) likelihood.  This is beyond the existing \(\Delta F=2\)
and \(\mu\to e\gamma\) surfaces and should target \texttt{quarkConstraints/modern/}
only if the PI later approves implementation.

\subsection*{Key references}
Process-local reference keys before bibliography consolidation:
\texttt{hflav\_ckm2025\_rdrds},
\texttt{hflav\_ckm2025\_logfile},
\texttt{babar2012\_arxiv1205\_5442},
\texttt{babar2013\_arxiv1303\_0571},
\texttt{belle2015\_arxiv1507\_03233},
\texttt{belle2017\_arxiv1612\_00529},
\texttt{belle2018\_arxiv1709\_00129},
\texttt{belle2019\_arxiv1910\_05864},
\texttt{lhcb2023\_arxiv2302\_02886},
\texttt{lhcb2023\_arxiv2305\_01463},
\texttt{lhcb2025\_arxiv2406\_03387},
\texttt{belleii2024\_arxiv2401\_02840},
\texttt{belleii2025\_arxiv2504\_11220},
\texttt{flag2024\_arxiv2411\_04268}, and
\texttt{cfw2008\_arxiv0804\_1954}.
